% GENERATED FILE. Edit canonical lessons, then run npm run generate:books.
% canonical-source-hash: fnv1a64:7f77a3c28173030b
% canonical-lessons: GE-C533-kamin, GE-C533-kommode, GE-C533-schublade, GE-C533-bettdecke, GE-C533-gemuetlichkeit, GE-R533-review-530-533, GE-R533-review-final-533

\chapter{Ein gemütliches Zuhause}
\label{ch:ge-ein-gemutliches-zuhause}
\hlchaptermodality{533}

\begin{chapteropening}
\textbf{By the end of this chapter:} \emph{I can describe a cosy home.}

Complete the last lesson of chapter 533: I can describe a cosy home.
\end{chapteropening}

\section[der Kamin]{der Kamin — fireplace; chimney}

\label{lesson:GE-C533-kamin}



\begin{quote}
Before the new one: say the German for caretaker, then the German for stairwell.
\end{quote}

\subsection*{You'll want to know: der Kamin}
\textbf{der Kamin} — \textquotedblleft{}fireplace; chimney\textquotedblright{}.

\begin{quote}
Im Winter sitzen wir abends vor dem Kamin.
\end{quote}

\begin{quote}
\emph{In winter we sit in front of the fire in the evenings.}
\end{quote}

Plural: \emph{die Kamine}. In southern Germany and Austria also the chimney.

\begin{cousinweb}
\emph{Kamin} is inherited from Middle High German \emph{kamīn} and Old High German \emph{kemīn}, an early borrowing from Latin \emph{camīnus} \textquotedblleft{}furnace\textquotedblright{}, from Greek \emph{káminos} (EtymDB).
\end{cousinweb}

\begin{grammarlens}[title={What you've built}]
One more word for telling, explaining and describing at B1.
\end{grammarlens}

\subsection*{Your turn}
\begin{itemize}
\raggedright
  \item \emph{Say it:} \emph{der Kamin}
  \item \emph{Say it:} \emph{der Kamin}, once more
  \item \emph{Say it:} a sentence of your own with \emph{der Kamin}
  \item \emph{Recall:} say the German for caretaker, then the German for stairwell, then say \emph{der Kamin} again
\end{itemize}

\begin{tcolorbox}[breakable,colback=teal!4,colframe=teal!35!black,title={Before you move on}]
What does der Kamin mean? (\textquotedblleft{}Fireplace; chimney\textquotedblright{}.) Say it once more.
\end{tcolorbox}
\section[die Kommode]{die Kommode — chest of drawers}

\label{lesson:GE-C533-kommode}



\begin{quote}
Before the new one: say the German for stairwell, then the German for fireplace.
\end{quote}

\subsection*{You'll want to know: die Kommode}
\textbf{die Kommode} — \textquotedblleft{}chest of drawers\textquotedblright{}.

\begin{quote}
Die Socken liegen in der Kommode im Schlafzimmer.
\end{quote}

\begin{quote}
\emph{The socks are in the chest of drawers in the bedroom.}
\end{quote}

Plural: \emph{die Kommoden}.

\begin{cousinweb}
\emph{Kommode} was borrowed from French \emph{commode} \textquotedblleft{}chest of drawers\textquotedblright{} (EtymDB), from the adjective \emph{commode} \textquotedblleft{}handy\textquotedblright{}, from Latin \emph{commodus} \textquotedblleft{}convenient\textquotedblright{}.
\end{cousinweb}

\begin{grammarlens}[title={What you've built}]
One more word for telling, explaining and describing at B1.
\end{grammarlens}

\subsection*{Your turn}
\begin{itemize}
\raggedright
  \item \emph{Say it:} \emph{die Kommode}
  \item \emph{Say it:} \emph{die Kommode}, once more
  \item \emph{Say it:} a sentence of your own with \emph{die Kommode}
  \item \emph{Recall:} say the German for stairwell, then the German for fireplace, then say \emph{die Kommode} again
\end{itemize}

\begin{tcolorbox}[breakable,colback=teal!4,colframe=teal!35!black,title={Before you move on}]
What does die Kommode mean? (\textquotedblleft{}Chest of drawers\textquotedblright{}.) Say it once more.
\end{tcolorbox}
\section[die Schublade]{die Schublade — drawer}

\label{lesson:GE-C533-schublade}



\begin{quote}
Before the new one: say the German for fireplace, then the German for chest of drawers.
\end{quote}

\subsection*{You'll want to know: die Schublade}
\textbf{die Schublade} — \textquotedblleft{}drawer\textquotedblright{}.

\begin{quote}
Die Schere ist in der obersten Schublade.
\end{quote}

\begin{quote}
\emph{The scissors are in the top drawer.}
\end{quote}

Plural: \emph{die Schubladen}. In Austria also \emph{die Lade}.

\begin{cousinweb}
\emph{Schublade} is built on \emph{schieben} \textquotedblleft{}to push\textquotedblright{}, from Proto-Germanic \emph{skeubaną}, and \emph{Lade} \textquotedblleft{}box, drawer\textquotedblright{}, inherited from Middle High German \emph{lade} (EtymDB).
\end{cousinweb}

\begin{grammarlens}[title={What you've built}]
One more word for telling, explaining and describing at B1.
\end{grammarlens}

\subsection*{Your turn}
\begin{itemize}
\raggedright
  \item \emph{Say it:} \emph{die Schublade}
  \item \emph{Say it:} \emph{die Schublade}, once more
  \item \emph{Say it:} a sentence of your own with \emph{die Schublade}
  \item \emph{Recall:} say the German for fireplace, then the German for chest of drawers, then say \emph{die Schublade} again
\end{itemize}

\begin{tcolorbox}[breakable,colback=teal!4,colframe=teal!35!black,title={Before you move on}]
What does die Schublade mean? (\textquotedblleft{}Drawer\textquotedblright{}.) Say it once more.
\end{tcolorbox}
\section[die Bettdecke]{die Bettdecke — duvet, quilt}

\label{lesson:GE-C533-bettdecke}



\begin{quote}
Before the new one: say the German for chest of drawers, then the German for drawer.
\end{quote}

\subsection*{You'll want to know: die Bettdecke}
\textbf{die Bettdecke} — \textquotedblleft{}duvet, quilt\textquotedblright{}.

\begin{quote}
Unter der dicken Bettdecke ist es schön warm.
\end{quote}

\begin{quote}
\emph{It's lovely and warm under the thick duvet.}
\end{quote}

Plural: \emph{die Bettdecken}.

\begin{cousinweb}
\emph{Bettdecke} is built on \emph{Bett} \textquotedblleft{}bed\textquotedblright{}, from Proto-Germanic \emph{badją}, and \emph{Decke} \textquotedblleft{}cover\textquotedblright{}, inherited from Middle High German \emph{decke} (EtymDB).
\end{cousinweb}

\begin{grammarlens}[title={What you've built}]
One more word for telling, explaining and describing at B1.
\end{grammarlens}

\subsection*{Your turn}
\begin{itemize}
\raggedright
  \item \emph{Say it:} \emph{die Bettdecke}
  \item \emph{Say it:} \emph{die Bettdecke}, once more
  \item \emph{Say it:} a sentence of your own with \emph{die Bettdecke}
  \item \emph{Recall:} say the German for chest of drawers, then the German for drawer, then say \emph{die Bettdecke} again
\end{itemize}

\begin{tcolorbox}[breakable,colback=teal!4,colframe=teal!35!black,title={Before you move on}]
What does die Bettdecke mean? (\textquotedblleft{}Duvet, quilt\textquotedblright{}.) Say it once more.
\end{tcolorbox}
\section[die Gemütlichkeit]{die Gemütlichkeit — cosiness}

\label{lesson:GE-C533-gemuetlichkeit}



\begin{quote}
Before the new one: say the German for drawer, then the German for duvet.
\end{quote}

\subsection*{You'll want to know: die Gemütlichkeit}
\textbf{die Gemütlichkeit} — \textquotedblleft{}cosiness\textquotedblright{}.

\begin{quote}
Kerzen und leise Musik sorgen für Gemütlichkeit.
\end{quote}

\begin{quote}
\emph{Candles and soft music make it cosy.}
\end{quote}

No plural. Built on \emph{gemütlich} — cosy, relaxed.

\begin{cousinweb}
\emph{Gemütlichkeit} is built on \emph{gemütlich} \textquotedblleft{}cosy\textquotedblright{}, inherited from Middle High German \emph{gemüetlich}, built on \emph{gemüete} \textquotedblleft{}mind, mood\textquotedblright{} (EtymDB), with \emph{-keit}.
\end{cousinweb}

\begin{grammarlens}[title={What you've built}]
One more word for telling, explaining and describing at B1.
\end{grammarlens}

\subsection*{Your turn}
\begin{itemize}
\raggedright
  \item \emph{Say it:} \emph{die Gemütlichkeit}
  \item \emph{Say it:} \emph{die Gemütlichkeit}, once more
  \item \emph{Say it:} a sentence of your own with \emph{die Gemütlichkeit}
  \item \emph{Recall:} say the German for drawer, then the German for duvet, then say \emph{die Gemütlichkeit} again
\end{itemize}

\begin{tcolorbox}[breakable,colback=teal!4,colframe=teal!35!black,title={Before you move on}]
What does die Gemütlichkeit mean? (\textquotedblleft{}Cosiness\textquotedblright{}.) Say it once more.
\end{tcolorbox}
\section[(dialogue)]{Freizeit und Zuhause}

\label{lesson:GE-R533-review-530-533}



\begin{quote}
Say the German for to comb and to gargle.
\end{quote}

\subsection*{Your turn}
Say these aloud:

\begin{itemize}
\raggedright
  \item to comb, to gargle, to put on make-up, bathrobe, pyjamas
  \item to warm up, microwave, kettle, to wash up, to dry
  \item shared flat, flatmate, house rules, caretaker, stairwell
  \item fireplace, chest of drawers, drawer, duvet, cosiness
\end{itemize}

\begin{tcolorbox}[breakable,colback=teal!4,colframe=teal!35!black,title={Before you move on}]
To comb? (\textbf{kämmen}.) To gargle? (\textbf{gurgeln}.) To put on make-up? (\textbf{schminken}.) Bathrobe? (\textbf{der Bademantel}.) Pyjamas? (\textbf{der Schlafanzug}.) To warm up? (\textbf{aufwärmen}.) Microwave? (\textbf{die Mikrowelle}.) Kettle? (\textbf{der Wasserkocher}.) To wash up? (\textbf{abwaschen}.) To dry? (\textbf{abtrocknen}.) Shared flat? (\textbf{die Wohngemeinschaft}.) Flatmate? (\textbf{der Mitbewohner}.) House rules? (\textbf{die Hausordnung}.) Caretaker? (\textbf{der Hausmeister}.) Stairwell? (\textbf{das Treppenhaus}.) Fireplace? (\textbf{der Kamin}.) Chest of drawers? (\textbf{die Kommode}.) Drawer? (\textbf{die Schublade}.) Duvet? (\textbf{die Bettdecke}.) Cosiness? (\textbf{die Gemütlichkeit}.)
\end{tcolorbox}
\section[(dialogue)]{Ein gemütliches Zuhause — once more}

\label{lesson:GE-R533-review-final-533}



\begin{quote}
Say the German for fireplace and chest of drawers.
\end{quote}

\subsection*{Your turn}
\begin{itemize}
\raggedright
  \item \emph{Say it:} fireplace, chest of drawers, drawer, duvet, cosiness
\end{itemize}

\begin{tcolorbox}[breakable,colback=teal!4,colframe=teal!35!black,title={Before you move on}]
Fireplace? (\textbf{der Kamin}.) Chest of drawers? (\textbf{die Kommode}.) Drawer? (\textbf{die Schublade}.) Duvet? (\textbf{die Bettdecke}.) Cosiness? (\textbf{die Gemütlichkeit}.)
\end{tcolorbox}
